\documentclass[conference]{IEEEtran}
\usepackage{cite}
\usepackage{pdflscape}  % Untuk orientasi halaman landscape
\usepackage{tabularx}   % Untuk lebar tabel otomatis
\usepackage{booktabs}   % Garis tabel standar akademis
\usepackage[table]{xcolor} % Untuk highlight baris

\begin{document}

% ==========================================
% TABEL REVIEW LITERATURE (SINGLE COLUMN LANDSCAPE)
% ==========================================
\clearpage
\onecolumn % Beralih sementara ke mode 1 kolom agar tabel melebar penuh
\pagestyle{empty} % Menghilangkan header/footer pada halaman tabel

\begin{landscape}
\vspace*{-1.5cm} % Menggeser posisi tabel sedikit ke atas
\begin{table}[p]
\centering
\caption{Ringkasan dan Perbandingan Literatur Survei NIDS dengan Penelitian yang Diusulkan}
\label{tab:literature_review}
\small
\renewcommand{\arraystretch}{1.35} % Mengatur jarak antar baris

% Memaksa tabel melebar memenuhi lebar kertas landscape
\makebox[\linewidth][c]{%
  \begin{tabularx}{1.15\textwidth}{>{\centering\arraybackslash}p{0.8cm} >{\centering\arraybackslash}p{1.2cm} >{\raggedright\arraybackslash}p{3.0cm} >{\raggedright\arraybackslash}X >{\raggedright\arraybackslash}p{3.0cm} >{\centering\arraybackslash}p{1.5cm} >{\centering\arraybackslash}p{2.2cm} >{\centering\arraybackslash}p{1.8cm} >{\centering\arraybackslash}p{1.5cm}}
  \toprule
  \textbf{No} & \textbf{Tahun} & \textbf{Peneliti} & \textbf{Metode} & \textbf{Dataset} & \textbf{Cross Dataset} & \textbf{Arsitektur} & \textbf{Adversarial} & \textbf{Real Trafik} \\

  \midrule
  \rowcolor[gray]{0.9} % Highlight baris kontribusi utama
  1 & \textbf{2025} 
  & \textbf{Penulis}
  & \textbf{XGBoost, \textit{Adversarial Training}, \textit{Saliency Map}} 
  & \textbf{CSE-CIC-IDS 2018} 
  & \textbf{Tidak}
  & \textbf{Centralized} 
  & \textbf{Ya}
  & \textbf{Tidak} \\

  \bottomrule
  \end{tabularx}%
}

\vspace{0.5cm}
\begin{minipage}{0.95\linewidth}
\hspace{0.5cm}\textbf{Referensi:}
\begin{enumerate}
    \item Alauthman M, Aslam N, Al-Qerem A, Aldweesh A, Sureephong P. Generative adversarial networks for intrusion detection systems: A comprehensive survey of applications, challenges, and research directions. \textit{Arabian Journal for Science and Engineering}. 2026 Jan;51(1):179--203.
    \item Alauthman M, Aslam N, Al-Qerem A, Aldweesh A, Sureephong P. Generative adversarial networks for intrusion detection systems: A comprehensive survey of applications, challenges, and research directions. \textit{Arabian Journal for Science and Engineering}. 2026 Jan;51(1):179--203.
\end{enumerate}
\end{minipage}

\end{table}
\end{landscape}

\clearpage
\twocolumn % Mengembalikan layout dokumen ke mode dua kolom standar IEEE untuk halaman berikutnya
\pagestyle{plain}

\end{document}